% 本文件由 LaTeX 排版 Agent 根据有效 translation.json 自动生成。
% author=Clement of Alexandria; work_id=0207; title=谁是那将得救的富人？

\chapter{谁是那将得救的富人？}

% structure: p0001 source=h1 target=omitted reason=page-title-already-rendered-by-parent

一、那些对富人献上赞美辞的人，在我看来，不仅被公正地视为谄媚者和卑贱之人——因为他们虚伪地硬说令人不悦的事使他们愉快——而且也是不虔敬和奸诈的。不虔敬，是因为他们既忽略赞美和光荣天主——唯有祂是完美和良善的，\Quote{万物都出于祂，借着祂，归于祂}\BibleRef{罗 11:36}——却将神圣的尊荣赐予那些沉溺于可憎可恶生活的人，而最要紧的是，这些人因此要受天主的审判；奸诈，是因为财富本身足以使拥有者骄傲并败坏他们的灵魂，使他们偏离救恩之路，而这些人却以过分赞美的快感膨胀富人的心，使他们对财富以外的一切全然鄙视，因为他们因财富而受钦佩；这正如俗语所说，火上浇油，在骄傲上堆砌骄傲，给财富这本来就沉重的负担又加上自负，而这负担本应减轻并除去，因为它是一种危险致命的疾病。因为那自高自大的人，随之而来的便是降卑和堕落，正如圣言所教导的。在我看来，远胜于卑贱地谄媚富人并因其恶行赞美他们，乃是以各种可能的方式帮助他们获得救恩；向天主祈求这事，祂必确实而甘甜地将这些赐给祂的子女；如此，借着救主的恩宠，医治他们的灵魂，光照他们，引领他们达到真理的掌握；谁获得这真理并在善工上卓越，必将赢得永生的奖赏。那坚持到生命末日的祈祷，需要坚强而平静的灵魂；而生活的实践，则需要良善正直的心志，向着救主的一切诫命伸展。\par

二、或许救恩对富人比对穷人显得更困难的原因，不止一个而是多个。因为有些人仅仅听到救主的话，且是随意地听到，说\Quote{骆驼穿过针孔，比富人进入天国还容易}\BibleRef{玛 19:24}，便对自己绝望，认为自己注定不得生存，将一切交给世界，紧抓着现世生命，仿佛只有它是留给他们的，从而更加偏离通往来世生命的道路，不再询问主和师傅究竟称谁为富人，也不问那在人不可能的事，在天主如何成为可能。但另一些人正确而充分地理解这话，却对引向救恩的善工不够重视，没有为达到他们所希望的作必要的预备。我所说的这两点，都是针对那些已经认识救主大能和祂光荣救恩的富人。至于那些不认识真理的人，我几乎不关心。\par

三、那些因爱真理和爱弟兄而行动的人，既不粗鲁地傲慢对待被召的富人，也不因自己贪婪的目的而谄媚他们，必须首先用话语解除富人无根据的绝望，并借着对主圣言的恰当解释，表明如果他们服从诫命，天国的产业并非完全与他们无分；然后要告诫他们，他们怀有莫名的恐惧，主乐意接纳他们，只要他们愿意；此外，还要展示并教导他们如何以及借着哪些行为和态度赢得所希望的事物，因为这既非遥不可及，也非无需努力就可得着；相反，正如运动员的情况——把小而易朽的事与伟大不朽的事相比——那拥有世俗财富的人应当认为这在于他自己。因为在运动员中，有人因绝望于能获胜并赢得冠冕，而不报名参赛；另有人心怀得胜的希望，却不接受相应的训练、饮食和锻炼，结果未能获得冠冕，期望落空。同样，那被赋予世俗财富的人，只要他是信徒并默观天主爱人的伟大，就不应一开始就宣称自己被排除在救主的竞赛场之外；另一方面，也不应期望不经过奋斗和努力就能抓住不朽的荣冠，继续不训练、不参赛。而要前往，将圣言当作自己的教练，将基督当作竞赛的主席；以主的新约作为他规定的饮食，以诫命作为锻炼，以美好的德性——爱、信德、望德、真理的知识、温和、谦卑、怜悯、庄重——作为优雅和装饰；这样，当最后的号角发出起跑和离世的信号，如同离开生命的竞技场时，他就能以良好的良心，在颁发奖赏的审判者面前得胜，显然配得那至高的家乡，他将带着荣冠和天使的欢呼回到那里。\par

IV. 愿救主因而恩准我们，从这一点开始这个主题，就希望本身和达到希望的方法，向弟兄们提供真实、适宜且救人的教诲。祂确实赐予祈求者，教导询问者，驱散无知，消除绝望，再次引入关于富人的相同话语，这些话语自身就是诠释者和无误的阐释者。因为没有什么比再次聆听那些直到如今在福音中令你们困扰的同样陈述更好的了，你们未经审视地听，由于幼稚而错误地听：\BibleQuote{祂在路上前行时，有一个人跑来跪在祂面前，说：\InnerQuote{善师，我该做什么善事才能获得永生？}耶稣对他说：\InnerQuote{你为什么称我为善？除了一个，就是天主，没有善的。诫命你都知道：不可奸淫，不可杀人，不可偷盗，不可作假见证，不可欺诈，要孝敬你的父母。}他回答说：\InnerQuote{老师，这一切我自幼就都遵守了。}耶稣定睛看他，就喜爱他，对他说：\InnerQuote{你若愿意成全，你还缺少一样：去，变卖你所有的一切，施舍给穷人，你必有宝藏在天上；然后来，跟随我！}他听了这话，面带愁容，忧忧郁郁地走了，因为他有许多产业。耶稣环视左右，对门徒说：\InnerQuote{有钱财的人进天主的国是多么难啊！}门徒就惊奇祂的话。耶稣又对他们说：\InnerQuote{孩子们，仗恃钱财的人进天主的国是多么难啊！骆驼穿过针眼，比富人进天主的国还容易。}他们越发惊奇，彼此说：\InnerQuote{这样，谁还能得救？}耶稣注视他们说：\InnerQuote{在人不可能，在天主却不然，因为在天主，一切都是可能的。}伯多禄开始对祂说：\InnerQuote{看，我们舍弃了一切，跟随了祢。}耶稣回答说：\InnerQuote{我实在告诉你们：凡为了我和福音的缘故，舍弃了自己的父母、兄弟、财产的，没有不在今时就得百倍的田地、财产、房屋、兄弟，连迫害也在内，并在来世获得永生的。但有许多在先的，将要成为在后的；在后的，将要成为在先的。}}\par

V. 这些事记载在\BibleBook{马尔谷}{福音}中，其余各卷也相应一致；虽然各卷措辞或许略有差异，但都显示意思完全相同。\par

但我们深知，救主教导一切并非以单纯属人的方式，而是以神圣而奥妙的智慧教导祂自己的子民，所以我们不可按血肉之躯来听祂的话，而必须经过适当的探究和领悟，去寻求并学习其中隐藏的意义。因为即便是那些似乎已被主亲自向门徒简化了的事，也发现需要不少于——甚至更多——比以隐晦方式表达的事还要多的关注，这是由于其中超越性的智慧之丰盛。而那些被认为是被祂向那些在内的人——祂称为天国之子的人——所解释的事，仍然需要比那些似乎被简单表达、因而听者没有提出疑问的事更多的思考；但这些事关乎整个救恩计划，并要以可赞叹的超天心智来默观，我们不可仅用耳朵肤浅地接受，而要运用心智直达救主的精神和话语中未言明的含义。\par

VI. 因为我们的主和救主被问了一个最合适祂的问题——生命被问及生命，救主被问及救恩，导师被问及所教导的主要教义，真理被问及真实的不朽，圣言被问及圣父的言词，成全者被问及圆满的安息，不死者被问及确实的不朽。祂被问及那些祂为此降世、祂所灌输、祂所教导、祂所提供的事，以显示福音的本质，即永生的恩赐。因为祂作为天主，预见了所问之事以及各人将如何回答。因为谁能比先知中的先知、众先知之神的主更能如此呢？当被称呼为\Quote{善}，并从这第一个词开始取起点，祂以此开始了祂的教导，引领门徒归向天主，那善的、首要的、唯一的永生分施者，而圣子从祂领受了永生，并赐给我们。\par

VII. 因此，有关生活之教导中最伟大、最首要的一点，必须从一开始就植入灵魂中——认识永恒的天主，永恒之物的赐予者，并通过认识和领悟而拥有天主，祂是首先的、至高的、唯一的、善的。因为这是生命不变不移的泉源和支柱，即认识那真实存在的天主，祂赐予真实存在的事物，也就是永恒的事物，从祂那里，存在及其持续被赋予其他存有。因为不认识祂就是死亡；而认识祂、拥有祂、爱祂、肖似祂，才是唯一的生命。\par

八、那么，凡愿意活出真正生命的人，首先被命令要认识祂，\BibleQuote{除了圣子，没有人认识父}。\BibleRef{玛 11:27}其次，需要学习救主之后的伟大和恩宠的新颖；因为，按照宗徒的教导，\BibleQuote{法律是藉梅瑟传授的，恩宠和真理却是由耶稣基督而来的}。\BibleRef{若 1:17}藉忠信的仆人赐予的恩赐，与由真正的圣子赐予的，并不相等。如果梅瑟的法律足以赐予永生，那么救主亲自来临、为我们受苦，从诞生到十字架完成人性的历程，便是徒劳的；而那自幼就遵守了法律所有诫命的人，跪下来向另一人祈求不朽，也是徒劳的。因为他不仅成全了法律，而且从极幼年就开始这样做。一个没有产生过过错的老年，有什么伟大或显著卓越的呢？但若有人在少年的嬉戏和青春的热火中，显示出超过年龄的成熟判断，那才是令人钦佩和卓越的斗士，在心灵上早已白发苍苍。\par

然而，尽管如此，这个人深信自己在义德方面一无所缺，但在生命方面却完全匮乏。因此，他向那位唯一能赐予生命者祈求。关于法律，他满怀自信；但对天主子，他却以恳求的语气说话。他从信德转移到信德。在法律中如同危险地漂泊、停泊在危险的锚地，他转向救主以寻找港湾。\par

九、因此，耶稣并没有责备他没有完全遵守法律的一切，反而爱他，并喜悦地接纳他在所学之事上的服从；但他说，就永生而言，他并不成全，因为他没有成就那成全的事，他是法律的实行者，却在真生命上无所事事。那些事确实是好的。谁能否认呢？因为\BibleQuote{诫命是圣的}，作为某种以畏惧和预备性训练为方式的陶冶，带领人走向立法的顶峰和恩宠。\BibleRef{罗 7:12}但基督是\BibleQuote{法律的终向，使一切信者获得义德}；祂不是像奴仆那样使人成为奴隶，而是使人成为儿子、兄弟和共同承继者，他们承行父的旨意。\BibleRef{迦 3:24}\par

十、\BibleQuote{你若愿意是成全的}。\BibleRef{玛 19:21}因此他尚未成全。因为没有比成全更成全的了。而\Quote{若愿意}这句话神圣地显示了与祂交谈的灵魂的自决。因为选择取决于人，人是自由的；但恩赐却取决于天主，祂是主。祂赐给那些愿意、极其热切并祈求的人，好使他们的救恩成为他们自己的。因为天主不强迫（强迫与天主相悖），而是供给那些寻求的，赐予那些求问的，给那些敲门的开门。你若愿意，那么，若你真的愿意，而不欺骗自己，就获得你所缺乏的。你所缺乏的只有一件事——那唯一恒常的、善的、如今超越法律的、法律所不赐予、法律所不包含的、属于那活着之人的特权。那个从幼年就满足了法律所有要求、并以这卓越之事为荣的人，却无法以救主特别要求的这一件事来完成全部，以获得他所渴望的永生。但他不悦地离开了，因生命的诫命而烦恼，正是他为此祈求的生命。因为他并非如他所声称的那样真正渴望生命，而只是追求那善的选择的虚名。他能够忙于许多事；但那一件事，生命的工程，他却没有能力、也不愿、不能完成。这正是主对玛尔大所说的话，她忙于许多事，因服侍而分心烦恼；她责怪妹妹，因为妹妹放下服侍，坐在祂脚前专心学习：\BibleQuote{你为了许多事操心忙碌，但玛利亚选择了更好的一份，是不能从她夺去的。}\BibleRef{路 10:41}同样，祂也吩咐他放下忙碌的生活，专心于一位，并依附那位赐予永生的主的恩宠。\par

十一、那么，是什么说服他逃离，使他离开师傅、离弃恳求、希望和之前热切追求的生命呢？——\Quote{变卖你的财产}。这是什么意思？祂并非如一些人仓促理解的那样，命令他扔掉所拥有的财物、放弃财产；而是命令他从灵魂中驱逐关于财富的观念、对财富的兴奋和病态情感，以及那些忧虑——它们是存在的荆棘，窒息生命的种子。因为如果没有特殊目的——除非是为了生命——那么贫乏本身并不伟大或可贵。因为那些一无所有、赤贫、靠每日乞食为生、散布在街上的穷人，他们不认识天主和天主的义德，仅仅因为极端的匮乏和缺乏生计，甚至最小的物品也欠缺，他们难道就是最有福、最亲近天主，而且是永生唯一的拥有者吗？\par

放弃财富并将其施舍给穷人或有需要的人，并非新鲜事；因为在救主降临之前，许多人已经这样做了——有些人是为了获得学习的闲暇和追求一种死的智慧；另一些人则是为了空洞的名声和虚荣，如阿那克萨戈拉斯、德谟克利特和克拉特斯等人。\par

XII. 那么，为何命令那些从前没有拯救过人的事，当作新的、神圣的、唯一赋予生命的事呢？天主的圣子所启示和教导的新受造物，其独特之处是什么？不是他人已经做过的外在行为，而是它所指向的另一种更大、更神圣、更完善的事：从灵魂本身和心意中剥去情欲，从根上砍断并驱逐与心神相异的一切。因为这是信徒特有的教训，是配得上救主的教诲。因为从前轻视外在事物的人，虽然放弃并挥霍了他们的财产，但我相信，他们却加剧了灵魂的情欲。因为他们沉溺于傲慢、虚饰和虚荣，并藐视其余人类，仿佛自己做了超人的事。那么，救主怎会命令那些注定要永生的人做那些对祂所应许的生命有害和伤害的事呢？因为即使如此，一个人在摆脱财富的重担后，仍可能具有与生俱来且活生生的贪财之欲和渴望；他可能放弃了使用财富，但既缺乏又渴望他已花费的，便会因缺少陪伴和悔恨的存在而加倍忧伤。因为缺乏生活必需品的人不可能不心中烦恼，并在设法从某处提供这些东西时，被阻碍去追求更美好的事物。\par

XIII. 相反的情况更有益得多：一个人拥有足够的生活资料，既不为钱财所困，又能帮助那些需要帮助的人！因为如果没有人拥有任何东西，那么人中间还能留下什么施舍的余地呢？这教义岂不显然与主其他许多卓越的教导相悖和冲突吗？\BibleQuote{你们要藉那不义的钱财结交朋友，为在你们匮乏时，他们能收留你们到永远的帐幕里。}\BibleRef{路 16:9}\BibleQuote{该在天上为自己积蓄财宝，那里没有虫蛀，没有锈蚀，也没有贼挖洞偷窃。}\BibleRef{玛 6:19}一个人若先把自己所有的一切都丢弃了，怎能给饥饿的人食物，给口渴的人水喝，给赤身露体的人衣服穿，给无住处的人房屋住呢？若不这样做，主就以火和外面的黑暗来威胁。不，祂还吩咐匝凯和玛窦，那些富有的税吏，以好客款待祂。祂不命令他们放弃财产，而是施行公正，除去不公正的判断，祂接着说：\BibleQuote{今天救恩临到了这一家，因为他也是亚巴郎之子。}祂如此称赞对财产的使用，以至于吩咐，除了这一点之外，还要分施，给口渴的人喝，给饥饿的人吃，收留无家可归者，给赤身露体者衣服。但如果不用物质就不能满足这些需求，而祂又命令人们放弃物质，那么主岂不是在劝人给予和不给予同一件事，喂养和不喂养，收留和拒绝，分享和不分享吗？这是最不合理的事。\par

XIV. 因此，财富若也能裨益邻人，就不该被抛弃。因为它们是财产，因为被拥有；是善物，因为有用且由天主为人的使用而提供；它们就在我们手中，置于我们的权下，如同材料和工具，对那些懂得工具的人可用于善用。你若巧妙地使用它，它就是巧妙的；你若缺乏技巧，它就因你的技巧缺乏而受影响，它本身无可指摘。财富就是这样的工具。你能正确使用它吗？它服务于正义。有人错误使用它吗？它便相反地成为不义的仆役。因为其本性是服役，而不是统治。因此，那本身无善无恶、无可指摘的东西，不该受责备；但能善用或滥用它的能力，由于拥有自由选择，才应受责备。这就是人的心神和判断，它在处理所赋予的事物时具有自由和自决。所以，愿没有人摧毁财富，而应摧毁与财富的更好使用不相容的灵魂情欲。这样，他成为有德、善良的人，就能善用这些财富。因此，变卖一切所有的，应理解为指灵魂的情欲而言。\par

XV. 那么，我要这样说。既然有些事物在灵魂之内，有些在灵魂之外，若灵魂善用它们，它们也被视为善；若恶用，则为恶——那么，那位命令我们舍弃财产的主，是排斥那些除去之后情欲仍存的事物呢，还是排斥那些除去之后财富甚至成为有益的事物呢？因此，若一个人抛弃了世俗财富，却仍能富足于情欲，即使没有外在物质供其满足——因为意向产生自己的效果，并窒息理性、压制它、用内在的私欲燃烧它——那么，他囊中羞涩而情欲富足，对他毫无益处。因为他所抛弃的并非应当抛弃的，而是中性的；他剥夺了自己有益的东西，却因缺乏外在满足手段而点燃了内在邪恶的燃料。因此，我们必须舍弃那些有害的财产，而不是那些能够服务的财产——如果人懂得正确使用它们。凡以智慧、节制和虔敬管理的，就是有益的；有害的必须抛弃。但外在事物并不造成伤害。所以，主引入了对外在事物的使用，命令我们抛弃的不是生活资料，而是那些滥用它们的东西。这些就是灵魂的软弱和情欲。\par

第十六．这些人的财富存在对所有人都是致命的，失去它则是有益的。为此，我们当使灵魂成为纯洁——即贫穷而赤裸——然后要听救主这样说：\Quote{来，跟随我。}因为对心地纯洁的人，祂现在成了道路。但天主的恩宠不能进入不洁的灵魂。那在情欲上富足、被许多世俗情感折磨的灵魂是不洁的。因为那持有财产、金银、房屋，视之为天主的恩赐，并用它们服务于那为救恩而赐予它们的天主；他知道拥有它们更多是为了弟兄而非自己；他超越对它们的拥有，而不是他所拥有之物的奴隶；他不将它们携带在灵魂中，也不将他的生命束缚和限定在它们之内，而是始终致力于某种善良神圣的工作，即使他必然有时被剥夺了它们，也能以愉快的心情承受它们的失去如同承受它们的丰裕。这就是被主祝福、被称为神贫的人，是天国合适的继承人，而非一个不能富裕生活的人。\par

第十七．但那个把财富携带在灵魂中，心中藏着的不是天主圣神而是金银或土地，总是无休止地获取财产，不断觊觎更多，弯腰低头，被世界的罗网束缚，本是尘土且注定归于尘土——这样的人怎能渴望并思念天国呢？他心里装着的不是心而是地或金属，他必然置身于他所选择的事物之中。因为人的心在哪里，他的宝藏也在哪里。主承认有两种宝藏——好的：\BibleQuote{因为善人从他心里的善库中发出善来；}和坏的：\BibleQuote{恶人从恶库中发出恶来：因为心里所充满的，口里就说出来。}\BibleRef{玛 12:34}既然宝藏与祂不是同一回事，与我们也是如此：那在发现时带来意外大收获的宝藏，另有第二种，是无益的、不可取的、有害的恶得；同样，也有善的富足和恶的富足，因为我们知道富足和宝藏在本性上并非彼此分离。一种富足是应拥有和获取的，另一种则不应拥有，而应抛弃。\par

同样，神贫是有福的。因此玛窦又加上：\BibleQuote{神贫的人是有福的。}\BibleRef{玛 5:3} 如何？\Quote{以心神。} 又说：\BibleQuote{饥渴慕求天主义德的人是有福的。}\BibleRef{玛 5:6} 因此，相反的那类穷人是有祸的，他们无分于天主，更无分于人间的财产，且未尝过天主的义德。\par

第十八．因此（\Quote{有钱人难以进入天国}这句话）应以学术的方式理解，而非笨拙、粗俗或肉体地理解。因为若这样理解，救恩就不在于外在事物，无论它们是多还是少、小还是大、显赫还是卑微、受尊重还是被轻视；而是在于灵魂的德行，在于信德、望德、爱德、友爱情谊、知识、温良、谦逊和真理，其回报就是救恩。因为任何人活着或死亡，并不是因为身体的美丑。而是那贞洁地、并按照天主使用所赐予的身体的人，要活着；那摧毁天主圣殿的人，必被摧毁。丑陋的人可能放荡，英俊的人可能节制。身体的强壮与高大不能使人存活，任何肢体也不能使人死亡。而是使用它们的灵魂提供了每个原因。经上说：有人打你的脸时，要忍受；\BibleRef{玛 5:39} 强壮健康的人可以做到。反之，软弱的人可能因脾气执拗而犯罪。同样，穷困潦倒的人可能沉醉于情欲；而富足世俗之人可能节制、在放纵方面贫穷、可靠、明智、纯洁、受管教。\par

既然如此，首先且首要的是灵魂要活着，而围绕它生发的德行拯救，恶行杀害；那么显而易见：借着在那些事物上贫穷（借着那些事物的财富，人毁灭它），灵魂得救；借着在那些事物上富足（那些事物的财富败坏它），灵魂被杀害。我们不要再从别处寻找结局的原因，而要从灵魂的状态与倾向中寻找，就是关于对天主的服从与纯洁，以及关于违反诫命和积累邪恶的方面。\par

XIX. 那真正并正当地富有的人，是因德行而富有，并能圣善而忠信地使用任何财富的人；而虚假富有的人，是按肉体富有，将生命变为外在的财产，这财产是短暂易逝的，时而属于这人，时而属于那人，最终不属于任何人。同样，也有真正的穷人，和另一种冒牌而虚假称为穷人的。那神贫的人，这才是正确的；而那在今世意义上贫穷的人，则是另一回事。对于那在今世财物上贫穷、却在恶习上富有的人，那在神上不贫穷、在天主面前不富有的人，有话说：抛弃你灵魂中那些外来的财产，好使心灵纯洁，得以看见天主；这无非是说：进入天国。你如何才能抛弃它们？借着变卖。那么，你是否要用财物换取金钱，通过交换财富，将你的有形财产变为金钱？绝非如此。而是要在你灵魂中原本固有的、你渴望拯救的东西的位置上，引入其他使人神圣化并侍奉永生的财富，即符合天主命令的性情；为此，你将获得无尽的赏报、尊荣、救恩和永垂不朽。你就这样正当地变卖那些多余的、关闭天国大门的财产，以那些能拯救你的财富来交换它们。让前者为属肉体的穷人所拥有，他们缺乏后者。但你，因领受属神的财富作为交换，如今将在天堂拥有财宝。\par

XX. 那个富有且遵纪守法的人，因没有以比喻的方式理解这些事，也不明白同一个人可以既贫穷又富有，既拥有财富又未拥有，既使用世界又未使用世界，便忧愁沮丧地离开了，抛弃了那种他只能渴望却无法达到的生命状态，为自己制造了不可能之事。因为灵魂不被围绕可观财富的奢华和迷人诱惑所引诱和毁灭，是困难的；但即使被财富包围，一个人仍有可能获得救恩，只要他使自己从物质财富中抽身出来，转向那以理智把握、由天主教导的财富，并学会正确地、恰当地使用无差别的世事，从而追求永生。连门徒们自己起初也惊慌诧异。他们为何听到这话如此惊慌？难道他们自己拥有大量财富吗？不，他们早已舍弃了他们的网、钩子和小船，那是他们仅有的财产。那么，他们为何惊慌地说：\Quote{谁能得救？} 他们作为门徒，已好好听懂了主以比喻和隐晦方式所说的话，并领会了言语的深度。因为他们因缺少财富而对得救充满希望。但当他们意识到自己尚未完全舍弃私欲偏情（因为他们还是新奉教者，刚被救主拣选不久）时，他们极其惊讶，对自己感到绝望，不亚于那个如此可怕地紧抓财富、宁要财富不要永生的富人。因此，门徒们完全有理由恐惧；如果那拥有财富和充满私欲偏情的人都是富人，并且两者都同样要被逐出天堂。因为救恩是属于纯洁和无私欲偏情的灵魂的。\par

XXI. 但主回答说：\Quote{因为在人不可能的，在天主却可能。} 这又充满极大的智慧。因为人靠自己努力和辛劳以求摆脱私欲偏情，一事无成。但若他明确表现出非常渴望和热切于此，则因天主力量的加入而达成。因为天主与甘愿的灵魂合作。但如果他们放弃了热忱，由天主赐予的神也会被抑制。因为强迫不愿者得救是施压者的行为；而拯救自愿者则是显示恩宠者的行为。天国也不属于贪睡和懒惰的人，\Quote{而是以猛力夺取的。} 因为只有这种暴力是可嘉的：强迫天主，从天主手中强取生命。而祂认识那些坚持不懈，甚至强取的人，便让步并赐予。因为天主喜爱在此类事上被战胜。\par

因此，有福的伯多禄——被选者、卓越者、门徒之首，救主曾为自己和祂单独为他付了殿税\BibleRef{玛 17:27}——一听到这话，便迅速抓住并领会了这句话。他怎么说？\Quote{看，我们舍弃了一切，跟随了祢。} 如果这里\Quote{一切}指他自己的财产，他是在夸耀自己舍弃了大约四枚铜钱，却忘记指明天国的赏报。但如果他们抛弃了我们刚才所说的那些，即旧有的心智财产和灵魂疾病，追随师傅的足迹，那么这就使他们与那些将要登记在天上的人联合。因为这样才算真正跟随救主：以无罪和祂的完美为目标，装饰并整理灵魂，如同在祂面前的一面镜子，在各个方面都相似地安排一切。\par

XXII. \Quote{耶稣回答说：\InnerQuote{我实在告诉你们：凡撇下自己的所有、父母、儿女和财富，为我和福音的缘故的，必要得到百倍的赏报。}} 但不要让这话困扰你们，也不要让另一处所说的更严厉的话困扰你们：\Quote{谁若不恼恨自己的父亲、母亲、儿女，甚至自己的性命，就不能做我的门徒。} \BibleRef{路 14:26} 因为和平的天主，祂也劝勉我们要爱仇敌，祂不会引人去恨恶并离弃至亲。但如果我们应当爱仇敌，那么按合理的推论，从仇敌往上，我们应当更爱那些亲族中最近的人。或者，如果我们应当恨恶血亲，那么推论告诉我们，我们更应当摒弃仇敌。这样，两种推论就会互相否定。但它们并不互相否定，也远非如此。因为出于同一情感和态度，并基于同一规则，一个爱仇敌的人可以恨恶父亲，因为他既不向仇敌报仇，也不敬重父亲超过基督。因为一句话铲除了仇恨和伤害，另一句话铲除了对亲属的羞怯感，如果这有害于救恩的话。因此，如果一个人的父亲、儿子或兄弟不敬神，成为信德的障碍和更高生活的阻碍，就让他不要与之友好或一致，而是因着属灵的敌意，断绝血肉的关系。\par

XXIII. 假设这是一场诉讼。让你的父亲想象自己出现在你面前，说：\Quote{我生养了你。跟随我，与我一同作恶，不要听从基督的律法。} 以及其他一个亵渎者、一个本性已死的人会说的话。\par

但另一方面，请听救主说：\Quote{我重生了你，你被世界生下来走向死亡。我释放了你，医治了你，救赎了你。我将向你显示良善圣父的面容。不要称呼地上的人为父。让死人埋葬死人；你跟随我。因为我要带你进入不可言喻、不可言说的福乐安息，那是眼所未见，耳所未闻，人心也未曾想到的；天使也渴望窥探，看看天主为圣徒和爱祂的儿女预备了什么美事。} \BibleRef{格前 2:9} 我是供养你的那一位，我把自己作为饼赐给你，凡尝过的就不再经历死亡，并且每日供应不死的饮品。我是超天课程的导师。为了你，我与死亡争斗，并偿还了你的死亡，那是你因从前的罪和不信天主所欠的。\par

听了双方这些考虑之后，你自己决断，并为自己的得救投上一票。若兄弟说同样的话，若子女、妻子或任何其他人如此说，让基督在你内得胜，胜过一切。因为祂为你而战。\par

XXIV. 你甚至可以对抗财富。说：\Quote{基督当然没有剥夺我的财产。主不会嫉妒。} 但你看见自己被他征服并推翻了吗？舍弃它，抛掉它，恨恶它，弃绝它，逃离它。\Quote{如果你的右眼使你跌倒，} 赶快 \Quote{把它剜出来。} \BibleRef{玛 5:29} 一人有眼而进入天国，比全身完整而投入烈火更好。无论是手、脚还是灵魂，恨恶它吧。因为如果它在此为基督的缘故被毁，它将在那边复活得生命。\par

XXV. 同样，下面的话也与此类似。\Quote{现在这个时候，没有田地、钱财、房屋和兄弟，连同迫害。} 因为主不是召叫那些一无所有、无家可归、无兄弟的人进入生命，祂也召叫了富人；但如我们上面所说，还有兄弟，如伯多禄与安德肋，雅各伯与若望——载伯德的儿子们，但他们彼此并同基督心意合一。而\Quote{连同迫害} 这一说法否定了拥有这一切中的任何一项。有一种来自外部的迫害，由人攻击信徒而起，或是出于仇恨、嫉妒、贪婪，或是出于魔鬼的煽动。但最痛苦的是内部的迫害，它出自人自己的灵魂，被不敬神的欲望、各种享乐、卑劣的盼望和毁灭性的梦所折磨；当它总是贪求更多，被兽性的爱所疯狂，被像刺棍和刺一样的激情所点燃时，它就被血覆盖，驱迫它走向疯狂的追求、对生命的绝望和对天主的蔑视。\par

这种来自内部的迫害更为严重和痛苦，它常与人同在，被迫害者无法逃脱；因为他随身处处携带着敌人。同样，来自外部的焚烧带来试炼，但来自内部的焚烧产生死亡。针对一个人的战争容易结束，但灵魂中的战争持续到死。\par

在这种迫害中，如果你拥有世俗的财富，如果你有血缘的兄弟和其他抵押品，就要放弃这一切引向邪恶的财富；为自己争取平安，使你自己摆脱持久的迫害；从这些转向福音；在一切之前选择你灵魂的救主、辩护者和\OriginalTerm{护慰者}{Paraclete}，生命之君。\Quote{因为那看得见的是暂时的，那看不见的是永远的。} \BibleRef{格后 4:18} 在现今的时候，事物是易逝和不稳固的，但在那将要来到的，是永生。\par

XXVI. \BibleQuote{在先的将成为在后的，在后的将成为在先的。} \BibleRef{谷 10:31} 这句话含义丰富，值得阐述，但目前无需探究；因为它不仅指富人，也显然指所有曾将自己交付于信仰的人。因此，暂且搁置一旁。但我认为，我们的论点已不亚于我们所承诺的：救主绝没有因财富本身以及拥有财产而排斥富人，也没有为他们关闭救恩之门；只要他们能够且愿意将自己的生活服从于天主的诫命，并将这些诫命置于转瞬即逝的事物之上，并且以稳定的目光注视主，如同那些期待好舵手示意的人，看他希望什么，命令什么，指示什么，向他的水手发出什么信号，从何处又将船驶向何方。因为一个人在归信之前，通过用心节省而积蓄了足够的生活所需，这有什么害处呢？或者，如果一个人一生下来，就由赐予他生命的天主将他安置在这样的人家中，处在富有、家产雄厚、财富显赫的人群中，这又有什么可指责的呢？因为如果由于无意中出生于富裕之家，一个人就被剥夺了生命，那么他反而受了创造他的天主的亏待，天主赐予他暂时的享受，却又剥夺了永生。如果财富是死亡的创造者和庇护者，那财富当初为何要从地上长出来呢？\par

但如果有人在财富之中能够摆脱其力量，保持谦逊之心，实行自我节制，唯独寻求天主，并呼吸天主、与天主同行，那么这样的穷人遵守了诫命，是自由、不受束缚、无病、不被财富所伤的。但若不然，\BibleQuote{骆驼穿过针眼，比富有的人进入天主的国还容易。} \BibleRef{谷 10:25}\par

那么，让骆驼在富人之前穿过狭窄的道路，或许意味着更高深的事理；这救主的奥秘将在\WorkTitle{首要原理与神学论}中学习。\par

XXVII. 好了，首先让我们呈现比喻中明显的一点，以及它为何被讲述。让这教导富足的人：他们不可忽视自己的得救，仿佛已经预定了灭亡，也不可将财富投入大海，或谴责它为生命的叛徒和敌人，而应学习如何使用财富并获得生命的方式和方法。因为既然没有人因恐惧自己是富人而必然丧亡，也没有人因相信并确信自己将得救而必然得救，那么，让我们看看救主为他们指定了何种希望，以及如何使意外之事得以肯定，所希望之事得以拥有。\par

因此，当有人问\BibleQuote{哪一条诫命是最大的？}时，主说：\BibleQuote{你要全心、全灵、全力爱上主你的天主；} \BibleRef{玛 22:36} 没有比这更大的诫命（祂说），而且理由极为充分；因为它命令关涉第一位、最伟大的那一位，即天主自己、我们的父，万物由祂造成、存在，得救者也归于祂。既然我们先被祂所爱，并领受了存在，那么将任何事物视为更古老或更卓越就是不敬的；我们只对这最大的恩惠回报这微小的感恩；并且无法想出任何其他报答天主的方式，因为祂一无所缺、完美无缺；而通过尽力爱父，我们获得了不朽。因为一个人越爱天主，他就越深入天主之内。\par

XXVIII. 其次，不大于此的诫命，祂说，是\BibleQuote{你要爱近人如你自己。} \BibleRef{玛 22:39} 因此，要爱天主胜过爱你自己。当问者问道：\BibleQuote{谁是我的近人？} \BibleRef{路 10:29} 祂并未像犹太人那样，指明血亲、同乡、皈依者、同样受割礼者或遵行同一法律的人。而是引入一个从耶路撒冷下到耶里哥的人，描述他被强盗打伤，半死不活地丢在路上，祭司经过却绕过，肋未人侧目而视，而那位被鄙视、被驱逐的撒玛黎雅人却怜悯了他；他不像那些人般随意走过，而是准备了遇险者所需之物，如油、绷带、牲口、给店主的部分现款和部分承诺。祂说：\BibleQuote{这三个人中，谁是那遭遇强盗者的近人呢？} 他回答：\BibleQuote{是怜悯他的那人。} \BibleRef{路 10:36} 因此，你也去照样做，因为爱德会萌发出善行。\par

XXIX. 在这两条诫命中，祂引入了爱；但按次序区分了它。在一条中，祂将爱的首位归于天主，将第二位分给我们的近人。除了救主自己，还能是谁呢？或者说，谁比祂更怜悯了我们？我们被黑暗的统治者几乎置于死地，带着许多创伤、恐惧、情欲、激情、痛苦、欺骗、享乐。这些创伤的唯一医生是耶稣，祂彻底根除了情欲——不像法律那样仅处理恶果、恶植物的果实，而是将斧头对准邪恶的根部。正是祂将酒（达味葡萄树的血）倒在我们受伤的灵魂上，带来了从圣父的慈悲中涌出的油，并丰盛地倾注。正是祂产生了健康与救恩的、不可解开的束带——爱德、信德、望德。正是祂使天使、掌权者、权势者为了极大的赏报而服事我们。因为他们也必因天主儿女荣耀的启示，从世界的虚空中被解救出来。因此，我们应当爱祂如同爱天主一样。而那承行祂的旨意、遵守祂诫命的人，就是爱基督耶稣。\BibleQuote{因为不是每一个对我说\InnerQuote{主啊，主啊}的人，都能进入天国；惟独那承行我父旨意的人。}\BibleRef{玛 7:21} \BibleQuote{你们为什么称呼我\InnerQuote{主啊，主啊}，却不照我所说的去做呢？}\BibleRef{路 6:46} \BibleQuote{但你们有福了，因为看见了和听见了义人和先知所未}（看见或听见的），\BibleRef{玛 13:16}，如果你们照我所说的去做。\par

XXX. 祂首先是爱基督的人；其次是爱和照顾那些信靠祂的人。因为凡为门徒所做的，主都视为为祂自己做的，并视之为全部。\BibleQuote{我父所祝福的，你们来吧！承受自创世以来为你们预备的天国。因为我饿了，你们给了我吃的；我渴了，你们给了我喝的；我作客，你们收留了我；我赤身露体，你们给了我穿的；我患病，你们看顾了我；我在监里，你们来探望了我。那时义人回答说：主啊！我们什么时候见了你饥饿而供养了你，或口渴而给了你喝的？我们什么时候见了你作客而收留了你，或赤身露体而给了你穿的？我们什么时候见了你患病，或在监里而来探望过你？君王便回答他们说：我实在告诉你们：凡你们对我这些最小兄弟中的一个所做的，就是对我做的。}\par

同样，在另一边，对那些没有做这些事的人，\BibleQuote{我实在告诉你们：凡你们没有给这些最小中的一个做的，便是没有给我做。} 在另一处，\BibleQuote{接纳你们的，就是接纳我；拒绝你们的，就是拒绝我。}\BibleRef{玛 10:40}\par

XXXI. 祂称这些人为儿女、儿子、小孩子、朋友和这里的小子，是指他们将来在天上的伟大。\BibleQuote{你们小心，不要轻视这些小子中的一个，因为他们的天使在天上，常见我在天之父的面。}\BibleRef{玛 18:10} 在另一处，\BibleQuote{你们小小的羊群，不要害怕！因为你们的父喜欢把天国赐给你们。}\BibleRef{路 12:32} 同样，祂说\BibleQuote{天国里最小的}——即祂自己的门徒——\BibleQuote{比妇女所生的若翰还大。}\BibleRef{玛 11:11} 又说，\BibleQuote{凡因义人或先知的名，接纳一个义人或先知的，必得义人或先知的赏报；凡因门徒的名，只给这些小子中的一个一杯凉水喝的，我实在告诉你们，他决失不了他的赏报。}\BibleRef{玛 10:41} 因此，这是唯一不失落的赏报。又说，\BibleQuote{你们要用不义的钱财结交朋友，这样，当你们穷困的时候，他们可以接你们到永远的帐幕里。}\BibleRef{路 16:9} 这表明，本质上，一个人在自己权力内所拥有的一切财产并非他自己的。而且从这不义中，被允许行出一件正义和救恩的事，即周济那些在圣父那里有永久居所的人。\par

看，首先，祂没有命令你被恳求或等待被催逼，而是你自己去寻找那些值得受益的人和救主相称的门徒。因此，宗徒的话也很出色：\BibleQuote{因为主喜爱乐意捐献的人。}\BibleRef{格后 9:7} 那人乐于施舍，毫不吝惜，播种以便也能这样收获，没有怨言、争论、后悔和分施，这是纯粹的善行。但比这更好的是主在另一处所说的话：\BibleQuote{求你的，就给他。}\BibleRef{路 6:30} 因为诚然，天主在施予中如此喜悦。这句话超越一切神性——不是等待被请求，而是亲自询问谁值得接受恩惠。\par

XXXII. 然后为慷慨指定这样的赏报——永恒的居所！啊，出色的交易！啊，神圣的货物！用金钱购买永生；通过施舍世上易逝之物，换取天上的永恒居所！航行到这个市场，如果你是明智的，富人啊！如果需要，绕过整个世界。不避危险和辛劳，为能在此购买天国。为何透明宝石和翡翠如此令你喜悦？还有那成为火的燃料、时间的玩物、地震的戏弄、暴君暴行的借口的房屋？渴望居住在天上，与天主一同为王。这个国度，一个效法天主的人会赐给你。在此处接受一点点，在那里，祂将使你永远与祂同住。求，使你可以得到；赶紧；努力；恐怕祂使你蒙羞。因为祂没有被命令接受，而是你被命令给予。主没有说：给予，或带来，或行善，或帮助，而是说：结交朋友。但朋友不是通过一次礼物，而是通过长期的亲密来证明自己。因为既不是一日之信德、爱德、望德，也不是一日之忍耐，而是\BibleQuote{那忍耐到底的，必得救。}\BibleRef{玛 10:22}\par

XXXIII. 那么人如何给出这些东西？因为我不但给予朋友，也给予朋友的朋友。谁是天主的朋友？你不要判断谁配得或谁不配得。因为你可能在你的意见上错误。在无知的不确定性中，与其提防那些不太好的人而错过遇见好人，不如为了值得的人而向不值得的人行善。因为若你为了节省并试图测试谁将配得地接受，你就有可能忽略一些天主所爱的人；其惩罚是永火的罚。但通过轮流供给所有需要的人，你必须必然在各方面找到一些在天主面前有力量拯救的人。\Quote{你们不要判断，免得你们受判断。因为你们用什么判断来衡量，也用什么来衡量你们；\BibleRef{玛 7:1} 好的尺度，压紧的，摇动的，满溢的，也要给你们。} 要向所有记载为天主门徒的人开放你的同情，不要轻蔑地看外貌，也不要对任何年龄阶段漫不经心。如果一个人看起来身无分文、衣衫褴褛、丑陋或虚弱，不要在心里为此烦恼并转开。这外形是从外给我们披上的，是我们进入这世界的机会，使我们能进入这共同学校。但里面住着隐藏的圣父，和祂的儿子，祂为我们死了并同我们一起复活了。\par

XXXIV. 这可见的外貌欺骗了死亡和魔鬼；因为内在的财富、美丽是它们所看不见的。它们为尸体而狂躁，轻视它为软弱，因为对内在的财富盲目；不知道我们带着\Quote{\BibleQuote{宝贝在瓦器里}}\BibleRef{格后 4:7}，由天主圣父的能力、天主圣子的血和圣神的露珠保护着。但是你这尝过真理并被算为配得伟大救赎的人，不要受骗。相反，与其他人不同，你要为自己收集一支没有武装、不好战、不流血、无激情、无玷污的军队：虔诚的老人、天主所爱的孤儿、以温良武装的寡妇、以爱装饰的男子。用你的钱财获得这样的守卫，为身体也为灵魂，因他们的缘故，一艘沉船得以浮起，仅由圣人们的祈祷掌舵；肆虐的疾病被制伏，因覆手而逃遁；强盗的袭击被解除武装，因虔诚的祈祷而被掠夺；恶魔的势力被粉碎，因强力的命令而羞愧地停止行动。\par

XXXV. 所有这些战士和守卫都是可信赖的。没有一个是闲散的，没有一个是无用的。一个能从天主那里为你求得宽恕，另一个在你生病时安慰你，另一个为你向万有之主哭泣和叹息同情，另一个教导一些有益于救恩的事，另一个以自信劝诫，另一个以仁慈劝告。所有都能真正地爱，无诡诈、无恐惧、无虚伪、无谄媚、无伪装。哦，可爱的灵魂的甜美服务！哦，有信心之心的有福思念！哦，惟独敬畏天主之人的真诚信德！哦，那些不能说谎之人的言语真理！哦，那些受托侍奉天主、说服天主、取悦天主、不触及你肉体的行为之美！要说话，但向住在你内的永恒之王说话。\par

XXXVI. 所有的信徒都是良善和像天主的，并配得上那包围他们如冠冕的名字。此外，还有另一些，是选民中的选民，并或多或少显著，因为他们像船只一样把自己从世界的波涛中拉上岸，并带到安全之地；他们不愿显得圣洁，若有人称他们为圣洁便羞愧；他们隐藏内心深处不可言喻的奥秘，不屑于让他们的高贵在世上被看见；圣言称他们为\Quote{\BibleQuote{世界的光，地上的盐}}\BibleRef{玛 5:13}。这就是种子，天主的肖像和模样，祂的真儿子和继承者，在这里如同寄居，由圣父至高管理和适宜安排所派遣，藉着祂，世界的可见和不可见之物被创造；有些是为了他们的服务，有些是为了他们的管教，有些是为了他们的教导；当种子留在这里时，万物都维系在一起；当它被聚集时，这些事物将很快消散。\par

XXXVII. 因为天主对爱的奥秘还有什么需要？那时你将要看到圣父的怀抱，只有天主的独生子曾揭示。天主本身就是爱；并出于对我们的爱而成为女性。在祂不可言喻的本质中，祂是圣父；在对我们的怜悯中，祂成为母亲。圣父因爱而成为女性：这伟大的证据就是祂从自身所生的那一位；爱所结的果实就是爱。\par

为此祂也降来。为此祂穿上了人。为此祂自愿服从人的经验，以便通过使自己达到我们所爱的软弱之尺度，祂也能相应地使我们达到祂自己力量的尺度。并且即将被献上，将自己作为赎价，祂为我们留下了一个新盟约：\Quote{我把我的爱赐给你们。} 这是什么，有多大？为了我们每一个人，祂给出了自己的生命——这等价于所有人。祂要求我们从彼此身上回报这爱。如果我们欠兄弟们生命，并且与救主立下了这样的相互契约，为什么我们要再囤积和封闭世间的财物，这些是贫穷的、非我们所有的、短暂的？我们岂要彼此封闭那不久将成为火的财产的东西？若望神圣而有力地说了：\Quote{\BibleQuote{不爱自己弟兄的，就是杀人者}}\BibleRef{若一 3:14}，是加音的种子，魔鬼的乳儿。他没有天主的怜悯。他没有更好事物的希望。他是贫瘠的；他是不结果子的；他不是永恒生命的超天葡萄树的枝条。他被砍下，等待永久的火。\par

XXXVIII. 但请学习那更卓越的道路，即保禄指明的获救之道。\Quote{爱不求己利，} \BibleRef{格前 13:5}却倾注于弟兄身上。她为他而焦急，为他而清醒地癫狂。\Quote{爱遮盖许多罪过。} \BibleRef{伯前 4:8} \Quote{圆满的爱驱逐恐惧。} \BibleRef{若一 4:18} \Quote{不夸张，不自大；不乐于不义，而乐于真理；凡事包容，凡事相信，凡事盼望，凡事忍耐。爱永不止息。先知之恩终必停止，语言之恩终必止息，治病之恩在地上终必消失。但如今常存的有信、望、爱这三样；其中最大的是爱。} 这很合理。因为信德在我们因看见而确信、因看见天主而消失时便会离去。望德在所盼望之事来临后便归于无有。然而爱德却达于圆满，且当那完美的事物赐下时，爱德越发增长。人若将爱德引入自己的灵魂，即使他生在罪过中，并做了许多禁止之事，藉着增加爱德并采纳纯净的悔改，他仍能挽回自己的过失。因为如果你明白了谁是那不得进入天堂的富人，以及他如何使用自己的财产，你就不应让这一事实使你陷入沮丧和绝望。\par

XXXIX. 如果有人能逃脱财富的过剩及其在生命道路上造成的困难，并能享受永恒的美善，却可能在领受印记与救赎之后，由于无知或非自愿的境况而陷入罪过或过犯，以至完全失足，那么这样的人便被天主完全弃绝。因为凡真心诚意、全心转向天主的人，门已敞开，那三倍喜乐的慈父迎接他真心悔改的儿子。真正的悔改是不再被相同的、曾使自己招致死亡的罪过所束缚，而是将其彻底从灵魂中根除。因为当这些罪过被铲除时，天主便重新在你内安住。经上记载：在天上与圣父及天使面前，为一个悔改的罪人，有极大的喜乐和欢庆。\BibleRef{路 15:10} 因此祂也呼喊：\Quote{我喜欢仁爱，而非祭献。} \BibleRef{欧 6:6} \Quote{我不愿罪人死，而愿他悔改。} \BibleRef{则 18:23} \Quote{你们的罪虽似朱红，必变成雪白；虽黑如墨，我必洗白如羊毛。} \BibleRef{依 1:18} 因为只有天主有能力赐予罪赦，不记过犯；主也命令我们每日宽恕悔改的弟兄。\BibleRef{玛 6:14} \Quote{你们纵然不善，尚且知道把好东西给你们的儿女，} \BibleRef{路 11:13} 何况仁慈之父，一切安慰的善父，祂极其怜悯、极其仁慈，长久忍耐，等待那些已悔改的人呢？而悔改实则是停止犯罪，不再往后看。\par

XL. 因此，天主赐予过去罪过的赦免；但未来的罪过，各人自己负责。悔改就是谴责过去的行为，恳求圣父使之忘却，唯有祂能藉由祂的怜悯撤销已做的事，并以圣神之露洗涤先前的罪过。\Quote{我凭你最后的状态审判你}，这也是结局在每一个案中大声宣告的话。所以即使有人在生活中行了极大的善事，却在临终时一头栽入邪恶，他先前的一切劳苦于他无益，因为在戏剧的终场他放弃了角色；而一个先前生活放荡堕落的人，若后来悔改，是可能在悔改之后的时间中克服长期以来的恶行的。但这需要极大的谨慎，正如长期患病的身躯需要调理和特别照料。盗贼，你想获得赦免吗？不要再偷盗。奸夫，不要再淫乱。娼妓，今后过贞洁的生活。抢夺者，归还，且归还比你拿的更多。假见证者，实践真理。发假誓者，不再起誓，并根除其余的私欲：忿怒、贪欲、忧愁、恐惧；为的是你在终了时能在此世先与反对者和好。不过，根深蒂固的情欲恐怕不可能一下子根除；但藉着天主的能力、人的转求、弟兄的帮助、真诚的悔改和不断的照料，它们会得以纠正。\par

XLI. 因此，你这个虚荣、有权势、富有的人，绝对有必要为自己安置一位属神的人，作你的训练者和指导者。敬畏他，哪怕只有一人；畏惧他，哪怕只有一人。将自己交付给聆听，哪怕只有一人坦率直言，使用严厉的言词，同时医治。因为对眼睛而言，并非总是放纵无度就好，有时流泪和疼痛反而对更大的健康有益。同样，对灵魂而言，没有比不间断的享乐更为有害的了。因为灵魂在融化中变得盲目，如果它不被坦率的言词所触动。当这个人发怒时，你要畏惧；当他叹息时，你要痛苦；当他平息怒气时，你要敬畏；当他恳求免罚时，你要抢先。让他为你度过许多不眠之夜，为你代祷于天主面前，以熟悉的连祷之魔力感动天父。因为当祂的儿女恳求怜悯时，祂不会拒绝。而他将为你纯洁地祈祷，被尊为天主的天使，不为你所忧，而是为你担忧。这才是真诚的悔改。\Quote{天主是轻慢不得的。} \BibleRef{迦 6:7} 祂也不理会虚浮的言语。因为唯独祂探查内心的骨髓和肺腑，垂听那在火中的人，倾听在鲸鱼腹中祈祷者；祂亲近所有相信的人，远离那不悔改的恶人。\par

XLII. 为使你更确信，如此真诚悔改后，你仍有确切的救恩希望，请听一则故事——虽似传说，实为叙事，由宗徒\Person{若望}流传并托付于记忆保管。当暴君死后，他从\Place{帕特摩岛}返回\Place{厄弗所}，应邀前往邻近的各族地区，在此处任命主教，在彼处整顿整个教会，在另一处按立那些被圣神所拣选的人。\par

他来到不远处的一座城（有人提及该城之名），并安顿了其他弟兄后，最后注视着被任命的主教，见一青年，体格强壮，容貌俊美，热情洋溢，便说道：\Quote{我在这教会面前，并以\Person{基督}为见证，郑重地将这青年托付于你。}对方应允并承诺一切后，他同样嘱咐并作证。随后他便起程前往\Place{厄弗所}。那位长老将托付给他的青年领回家中，养育、守护、珍爱，并最终为他付了圣洗。此后，他以为所加于青年身上的主的印记已是对他完全的保障，便放松了更严格的看管与守护。但青年过早获得自由后，一些同龄的游手好闲、放荡不羁、精于作恶的青年引诱了他。他们先以许多昂贵的宴席诱惑他；随后，夜间外出拦路抢劫时，也带他同行。接着他们竟敢共同干更大的事。青年逐渐习以为常，因天资刚强，一旦偏离正道，便如倔强有力的马，口衔嚼环，力冲深渊。他完全绝望于天主的救恩，不再思虑小事，既已丧亡，便决心与其余人同归于尽。他聚集他们，组成一伙强盗，成为凶残至极、嗜血成性、冷酷无情的匪首。\par

时光流逝，因出现某些需要，人们再次派人请\Person{若望}。他处理完他来此要办的其他事务后，说道：\Quote{来，主教，请将我和\Person{救主}当着你的教会——你所治理的，作为见证——托付给你的存款归还我们。}对方起初困惑，以为是无中生有的钱财指控，他并未收受；既不能相信未曾有的事，又不敢不信\Person{若望}。但当\Person{若望}说：\Quote{我要那青年，那弟兄的灵魂，}老人深深叹息，泪流满面，说：\Quote{他死了。}\Quote{如何死的？怎样死的？}\Quote{他死了，}他说，\Quote{向天主死了。因为他变得邪恶、放荡，终成强盗；如今他带领一伙同党占据了教堂前的山头。}宗徒撕裂衣服，捶头痛哭，说道：\Quote{我留下了一个多么好的弟兄灵魂的守卫！但给我备一匹马，让某人给我引路。}他当即骑马，径直从教堂出发。来到那地方，被强盗的哨兵截住；他既不逃跑也不求饶，却喊道：\Quote{我正是为此而来。带我去见你们的头目。}此时那首领正全副武装地等着。但他认出前来的\Person{若望}时，羞愧转身逃窜。\Person{若望}不顾年迈，全力追赶，喊道：\Quote{我儿，你为何逃避我——你手无寸铁、年迈的父亲？我儿，可怜我。不要怕；你仍有生命的希望。我愿在\Person{基督}面前为你负责。若有必要，我甘心为你死，如主为我们而死。我将为你舍弃性命。站住，相信吧；\Person{基督}派遣了我。}\par

他听见后，先低头站住；然后扔下武器，继而颤抖，痛哭流涕。老人走近时，他拥抱他，用自己的话尽力哭着辩护，只隐藏着右手。\Person{若望}向他保证，并起誓他必从\Person{救主}获得宽恕，恳求、跪拜，亲吻他那只已被悔改洁净的右手，将他带回教堂。然后，\Person{若望}以大量祈祷恳求，与他一同守斋克苦，用各种言语驯服他的心，直到最终——如他们所言——将他交还给教会，在他身上呈现了一个真正悔改的伟大榜样，一个重生的伟大记号，一个我们所盼望的复活的凯旋；当世界末日，天使们喜乐放光，歌唱着敞开天空，将那些真正悔改的人接入天上的居所；而\Person{救主}亲自走在前头迎接他们，欢迎他们；高举那无影、不止息的光；带领他们到\Person{圣父}的怀抱，进入永生，进入天国。\par

让人相信这些事，相信天主的门徒，相信作担保的天主，相信\WorkTitle{预言}、\WorkTitle{福音}、宗徒的言语；与之相合而生活，侧耳聆听，实践善行，他在离世时必看见所教导真理的结局与明证。因为在此世迎接悔改天使的人，在离开肉身之时必不后悔，当看见\Person{救主}带着荣耀和天军临近时，必不羞愧。他不惧怕火。\par

但如果有人选择继续在享乐中不断犯罪，且将此世的放纵看得比永生更重，并转离那赐予宽恕的救主，就不要再责怪天主、财富或他的堕落，而要责怪他自己的灵魂，因它是自愿灭亡的。但那位将目光转向救恩并渴求它、且以大胆和热切祈求赐予救恩的人，在天上的圣父必将赐予真正的净化与不变的生命。愿光荣、尊威、权能、永远的王权，藉祂的圣子\Person{耶稣基督}——生者与死者的主——并藉圣神，归于祂，从今世直到永远，世世代代，无穷无尽。阿们。\par

\SourceRecord{Translated by William Wilson.；Ante-Nicene Fathers；Vol. 2.；Edited by Alexander Roberts, James Donaldson, and A. Cleveland Coxe.；Buffalo, NY: Christian Literature Publishing Co.,；1885.；<http://www.newadvent.org/fathers/0207.htm>.}
